\section{Introduction}
\label{sec:intro}

We consider the problem of sampling the target posterior distribution with
density $\rho_\post (\cdot)$, for the parameter $\theta \in \R^{N_\theta}$,
given by
\begin{equation*}
    \rho_\post (\theta) \propto \exp (-V(\theta)), \qquad \theta \in \R^{N_\theta},
\end{equation*}
where $V(\theta) : \R^{N_\theta} \to \R$ is a known function while the
normalizing constant $Z = \int \exp (-V(\theta)) \dd \theta$ is intractable.
Markov chain Monte Carlo (MCMC) methods \cite{brooks2011handbook} aim to sample
from the target distribution $\rho_\post$ asymptotically, but can be
computationally expensive in high-dimensional or ill-conditioned settings
\cite{JMLR:v23:20-357}. Variational inference (VI) \cite{blei2017variational},
on the other hand, replaces exact sampling by optimization over a tractable
family of probability distributions on its Kullback--Leibler (KL) divergence to
the true target posterior, which is often more scalable than MCMC
\cite{JMLR:v14:hoffman13a}.

We focus on the Gaussian variational family
\begin{equation*}
    \PG := \left\{ \rho_a = \N(m,C) : a=(m,C)\in \Aspace := \R^{N_\theta} \times \SPD(N_\theta) \right\},
\end{equation*}
which is the simplest nontrivial VI family, where the Gaussian variational
inference problem is
\begin{equation*}
    a_\star = (m_\star,C_\star) \in \argmin_{a\in\Aspace} \calE(a),
    \qquad \calE(a) := \KL(\rho_a\,\|\,\rho_{\post}).
\end{equation*}

Recent theoretical work has studied Gaussian variational inference through
Wasserstein and Bures--Wasserstein geometry
\cite{lambert2022variational, diao2023forward}. This viewpoint connects Gaussian
variational inference with optimal transport and constrained Wasserstein
gradient flows, and yields a convergence theory for both the Wasserstein
gradient flow and its forward--backward discretizations. On the other side,
the Fisher--Rao gradient flow \cite{carrillo2026fisher} develops a complementary
information-geometric viewpoint \cite{amari2016information}. When the
Fisher--Rao gradient flow is restricted to the Gaussian family, we obtain the
natural gradient flow \cite{amari1998natural} for Gaussian variational
inference \cite{chen2023sampling},
\begin{equation}\label{ODEs:NG}
    \begin{aligned}
        \frac{\dd m_t}{\dd t} &= C_t\E_{\rho_{a_t}}[\nabla_\theta \log \rho_{\post}(\theta)], \\
        \frac{\dd C_t}{\dd t} &= C_t + C_t \E_{\rho_{a_t}}[\nabla_\theta \nabla_\theta \log \rho_{\post}(\theta)] C_t.
    \end{aligned}
\end{equation}
The Fisher--Rao metric restricted to the Gaussian family is affine invariant
\cite{chen2023sampling}: the flow \eqref{ODEs:NG} commutes with every invertible
affine change of variables $\theta \mapsto A\theta + b$, so that its behavior is
unchanged by an arbitrary linear reparameterization of the target. This is the
structural property that distinguishes it from the Bures--Wasserstein flow, and
it is why we expect natural gradient methods to be insensitive to the
conditioning of the target. It also determines which quantities can appear in a
sharp convergence rate. However,
while the natural gradient flow is useful in practice, its convergence theory
for Gaussian variational inference remains underexplored. We develop that theory
here, for the flow \eqref{ODEs:NG} and for two discretizations of it, in both
the deterministic and the stochastic setting.

The convergence of \eqref{ODEs:NG} and of its discretizations proceeds in three
stages,
\begin{equation}\label{eq:three-stage}
    \boxed{\;
    \text{covariance burn-in}
    \;+\;
    \text{global localization}
    \;+\;
    \text{local convergence},
    \;}
\end{equation}
each governed by a different mechanism and each carrying its own sharp rate. The
first stage is an artifact of the geometry rather than of the target. In the
covariance equation of \eqref{ODEs:NG} the growth of a small covariance is
multiplicative, and the mean equation carries the factor $C_t$, so an
underdispersed initialization suppresses the mean motion until the covariance
has reached the curvature scale of the target. The second stage is the global
descent proper, and it is controlled not by the curvature ratio
$\kappa = \beta/\alpha$ of the target but by the intrinsic ratio $\kappastar$
obtained after whitening by the optimizing covariance $C_\star$; the two agree
for a poorly conditioned Gaussian target only in the sense that $\kappastar = 1$
there, however ill-conditioned $\kappa$ may be. The third stage is local, and its
rate is governed neither by $\kappa$ nor by $\kappastar$ but by the spectral gap
of a self-adjoint operator built from the first- and second-Hermite components
of the whitened Hessian, which measures the non-Gaussian coupling between the
mean and covariance modes. The three mechanisms are logically independent: a
transport warm start removes the first without touching the other two.

The deterministic theory occupies \Cref{sec:global,sec:local} and gives matching
upper and lower bounds at every stage. We prove in \Cref{thm:glob-cont} an
affine-invariant continuous-time estimate that separates the initial covariance
from the intrinsic conditioning, and we show in \Cref{thm:sharp-cont} that its
$\kappastar$ dependence cannot be improved, by means of a two-dimensional
logarithmic-spiral potential. For the two discretizations, a Riemannian
retraction scheme and a KL/Bregman resolvent scheme, we prove in
\Cref{thm:glob-Riem,thm:glob-KL} global complexity bounds of a common form, and
we show in \Cref{thm:sharp-disc} that their factor $\kappa/\Delta t$ is sharp at
a fixed stepsize, by means of a one-dimensional bump-train potential that is
even, strongly convex, and uniformly Hessian-Lipschitz. The covariance burn-in,
and only the covariance burn-in, can be removed algorithmically: we give in
\Cref{thm:burnin-BW} a Bures--Wasserstein transport warm start and in
\Cref{thm:burnin-rescue} a single gradient--Hessian query that lands the
covariance on the curvature band, which for a Gaussian target recovers the exact
optimizer in one query by \Cref{thm:rescue-Gauss}.

The local theory keeps the nonlinear Gaussian covariance dynamics exactly and
perturbs only the non-Gaussian remainder. We prove in
\Cref{prop:spectral-sandwich} a two-sided spectral estimate for the linearized
generator, and we show in \Cref{thm:loc-sharp} that its $\betastar-1$ branch is
attained: a smooth convex-ridge family in dimension
$N_\theta = \Theta(\kappastar^2)$, with a dimension-free Hessian-Lipschitz
constant, has local gap $\Theta(\kappastar^{-1/2})$, which rules out every
dimension-free logarithmic local rate. The Gaussian-core decomposition gives in
\Cref{thm:loc-flow,thm:loc-disc} certified local regions that are exact for
Gaussian targets in every dimension and degrade continuously for general
strongly log-concave targets, together with an optimizer-free entry test in
\Cref{prop:entry-residual}. \Cref{thm:three-stage} assembles the three stages
into a single complexity statement for the flow and for both schemes.

The stochastic theory occupies \Cref{sec:stochastic} and reproduces the same
three stages under a joint Price/Hessian oracle with a sticking-the-landing
mean estimator \cite{roeder2017sticking}. Because the sampled Hessian obeys an
almost-sure spectral bound, the covariance bands hold pathwise with no clipping
and no stopping time, and the global theorems
\Cref{thm:stoch-glob-Riem,thm:stoch-glob-KL} therefore hold in plain expectation,
with explicit residual floors in two variance regimes. The local statement is
weaker than its deterministic counterpart: we prove in
\Cref{thm:stoch-loc-Riem,thm:stoch-loc-KL} a mean-square contraction on the
non-exit event together with a finite-horizon containment bound in
\Cref{prop:exit-prob}, and we do not claim infinite-horizon containment, since
the mean-score noise has unbounded support. The additive floors vanish for
Gaussian targets, and a decreasing stepsize removes them in general by
\Cref{thm:decreasing}. \Cref{cor:stoch-three-stage} is the stochastic mirror of
\Cref{thm:three-stage}.

Two further results mark the boundaries of the theory. Beyond log-concavity, the
unprojected schemes genuinely fail: a smooth double well with uniformly bounded
Hessian produces the covariance cascade of \Cref{prop:cascade-Riem} and the
resolvent pole of \Cref{prop:pole-KL}, while a covariance-constrained KL step
retains an $O(1/N)$ stationarity guarantee by \Cref{thm:clip-KL}. Finally, in
\Cref{sec:geometry} we classify, for $N_\theta \geq 2$, the complete
three-parameter family of affine-invariant Riemannian metrics on the Gaussian
manifold, compute in \Cref{thm:aff-modes} the exact rates at which each metric
preconditions the mean, trace, and traceless covariance modes on a Gaussian
target, and show in \Cref{thm:aff-balanced} that Fisher--Rao is the unique
member of the family, up to scale, that balances all three. This recovers the
metric with which we began.

\subsection*{Related work}

Natural gradients and information geometry provide the classical background for
the Fisher--Rao metric \cite{rao1945information, amari1998natural,
amari2016information}, and Gaussian and stochastic variational inference are
surveyed in \citet{blei2017variational} and \citet{JMLR:v14:hoffman13a}. The
Fisher--Rao gradient flow and its Gaussian restriction are studied in
\citet{carrillo2026fisher} and \citet{chen2023sampling}, the latter establishing
the affine invariance we exploit throughout. Our analysis differs from a generic
natural-gradient convergence argument in that it uses the exact Gaussian
covariance equation rather than a smoothness surrogate, and keeps every constant
in optimizer-whitened, affine-invariant form.

Wasserstein Gaussian variational methods supply a complementary geometry.
\citet{lambert2022variational} study variational inference through
Bures--Wasserstein gradient flows and analyze a stochastic scheme with
non-asymptotic guarantees, and \citet{diao2023forward} analyze a forward--backward
scheme with state-of-the-art rates and the first stationarity guarantees beyond
log-concavity. We use the Bures--Wasserstein flow in two limited roles, as a
warm start that removes the covariance burn-in and as a point of comparison
beyond log-concavity; we do not blend the two geometries into a
Wasserstein--Fisher--Rao method here. Affine-invariant metrics on the cone of
positive-definite matrices were classified in several forms by
\citet{thanwerdas2019affine, thanwerdas2023n}; our classification includes the
mean block of the Gaussian family and, for $N_\theta \geq 2$, uses it to expose
the exact modal rates.

The stochastic construction rests on the sticking-the-landing principle of
\citet{roeder2017sticking}, and natural-gradient variational inference with
nonconjugate models and constrained covariance parameterizations is studied in
\citet{sun2025natural}. Our stochastic results treat the joint geometry of the
mean and the full covariance, separate constant-step residual levels from
decreasing-step convergence, and keep local containment separate from local
contraction.

\subsection*{Organization and notation}

\Cref{sec:global} develops the global deterministic theory, its two
discretizations, the removal of the covariance burn-in, and the non-log-concave
boundary. \Cref{sec:local} gives the local spectral and nonlinear theory and the
three-stage complexity theorem. \Cref{sec:stochastic} treats the Price/Hessian
stochastic algorithms. \Cref{sec:geometry} classifies the affine-invariant
metrics, \Cref{sec:numerics} presents numerical illustrations, and
\Cref{sec:discussion} collects open problems. Longer proofs and the three lower
constructions are given in the appendices.

We write $N_\theta$ for the dimension, $\norm{\cdot}_\op$ and
$\norm{\cdot}_\Frob$ for the operator and Frobenius norms, and
$\DeltaE(a) = \calE(a) - \calE(a_\star)$ for the energy gap, with
$\DeltaE_0 = \DeltaE(a_0)$. The symbol $\kappa = \beta/\alpha$ always denotes the
curvature ratio in the original coordinates and $\kappastar$ its
optimizer-whitened counterpart, and similarly for $\alphastar$ and $\betastar$.
We write $\Delta t$ for the stepsize, $B$ for the batch size, and
$\logp x := \max\{\log x, 0\}$. Constants hidden by $O(\cdot)$ are numerical
unless their dependence is displayed.
